% ICLR 2026 submission.
%
% The style files in this directory are the OFFICIAL ones, downloaded from
% https://github.com/ICLR/Master-Template/tree/master/iclr2026 -- do not edit
% them. iclr2026_conference_ORIGINAL.tex is the upstream instructions file, kept
% for reference only; it is not part of this submission.
%
% Every table under tables/ is GENERATED. Do not edit them by hand:
%   python scripts/make_boba_paper_tables.py --analysis output-boba/analysis
% \todo{} marks everything not yet backed by a completed run.
%
%   cd paper && pdflatex main && bibtex main && pdflatex main && pdflatex main

\documentclass{article}
\usepackage{iclr2026_conference,times}

\usepackage{amsmath,amssymb}
\usepackage{booktabs}
\usepackage{graphicx}
\usepackage{xcolor}
\usepackage{hyperref}
\usepackage{url}

\newcommand{\todo}[1]{\textcolor{red}{[TODO: #1]}}
\newcommand{\optz}{\mathrm{opt}_z}
\newcommand{\frag}{\mathrm{frag}}

\title{What Noisy Feedback Costs Bayesian\\ Optimization, and the One Number\\ That Predicts It}

% Authors must not appear in the submitted version; keep \iclrfinalcopy
% commented out until camera-ready.
\author{\todo{authors} \\
Affiliation \\
Address \\
\texttt{email}
}

%\iclrfinalcopy % Uncomment for camera-ready ONLY, never for submission.
\begin{document}

\maketitle

\begin{abstract}
Bayesian optimization driven by human feedback is evaluated against regression
oracles fitted to archival ratings. That entangles the cost of noisy feedback
with the cost of the surrogate human being wrong, and because it is done on one
design space at a time it cannot say what generalises. We instead run four
human-plausible error processes on twenty analytic landscapes with published
optima, standardised so one unit of error means the same thing on each:
76{,}800 paired noisy-versus-clean comparisons over four magnitudes, two onsets
and twelve acquisition functions.
%
Bayesian optimization is more robust than the applied framing implies. Error
with the same standard deviation as the landscape itself, present from the very
first rating, costs $14.1\%$ of the improvement that was available; the same
error arriving after convergence costs $2.8\%$. Ordered by how much they move
that cost, the factors are error magnitude ($23\times$), landscape identity
($8\times$), onset ($5\times$), acquisition function ($2.8\times$) and the
\emph{kind} of error ($1.3\times$) --- the two levers practitioners pull are the
two smallest.
%
Almost all of the variation is captured by one statistic computable before any
optimization is run: the expected loss from a single noisy greedy choice,
$\frag(\sigma_e)$. Conditional on it, error magnitude has no residual predictive
power ($\beta = -0.07$, $p = 0.44$).
%
Two ablations say what the number means. Supplying the surrogate with the true
observation variance removes only $5\%$ of the measured cost, so this is an
information cost and not a mis-specified model; changing the incumbent
definition removes $15.6\%$, and $35$--$41\%$ for the probability-of-improvement
family, so published robustness rankings over acquisitions are rankings over
acquisition-and-incumbent pairs. Within-rater noise measured in three real
studies falls at $0.74$ to $3.35$ landscape standard deviations --- the range in
which we measure a $12$--$20\%$ loss.
%
Running the fitted-oracle design on the same grid puts the cost at roughly a
third of what the exact functions report, and disabling the oracle augmentation
--- worth up to $0.37$ held-out $R^2$ --- changes that by almost nothing. The gap
is a property of the design, not of how good the surrogate human is.
\end{abstract}

\section{Introduction}

Bayesian optimization is the standard tool for optimizing an expensive
black-box function from few evaluations \citep{jones1998efficient,
shahriari2016taking, frazier2018tutorial}, and it is increasingly pointed at
objectives that only a person can evaluate: interface parameters judged for trust or aesthetics, generated
content ranked for quality, robot behaviours rated for comfort
\citep{brochu2010tutorial, gonzalez2017preferential}. The optimizer's premise is
that each observation is informative about the objective. A person supplying
those observations is not a sensor. They rate inconsistently, they drift over a
session, they are systematically generous or harsh, and their errors are
correlated across consecutive judgements. The natural question --- how much does
that cost you --- is asked routinely and, we argue, answered in a way that
cannot support the answer.

The standard construction fits a regression model to archival ratings, treats it
as the objective, injects synthetic error on top of its predictions, and
measures the degradation. This has two problems that no amount of care in the
optimization code addresses. The first is entanglement: the fitted oracle is
itself wrong, often badly --- in the studies motivating this work its
cross-validated $R^2$ is $0.55$ at best and negative on one dataset --- so the
measured degradation mixes the cost of the injected error with the cost of the
surrogate's own mis-specification, and the reference optimum that regret is
measured against is a random-search estimate the optimizer can legitimately
beat. The second is generality: the experiment runs on one design space, so a
result such as ``acquisition $X$ is the most robust'' is a statement about one
landscape with no way to tell whether it is a statement about anything else.

We remove both by giving up the human entirely. On an analytic benchmark with a
published optimum the objective is exact, the regret is a real regret, and ---
the part that makes this more than a sanity check --- the geometry of the
landscape can be measured \emph{before} the optimization is run and used as a
predictor of what the error will cost. Running twenty such landscapes turns
``which acquisition is robust here'' into ``what property of a problem decides
whether feedback error hurts'', which is a question with a transferable answer.

Four findings follow. Bayesian optimization tolerates far more feedback error
than the applied framing suggests, and the cost is governed by \emph{when} the
error arrives about as much as by how large it is, while the \emph{kind} of
error --- random, systematic, drifting, autocorrelated --- barely matters.
Almost all of the variation is captured by a single cheap statistic of the
landscape, the expected loss from one noisy greedy pick, and conditional on it
the error magnitude carries no further information. The landscape matters
roughly three times more than the choice of acquisition function, and which
acquisition wins is itself mostly a property of the landscape. And a substantial
part of what looks like acquisition fragility is an artefact of how the
incumbent is defined, which is a free choice rather than a property of the
method.

\paragraph{Contributions.}
\begin{enumerate}
  \item A protocol for measuring the cost of feedback error with no fitted
    oracle in the loop: twenty analytic landscapes with verified optima on a
    common scale, four human-plausible error processes, and a model-free control
    whose excess regret is provably --- and measurably --- exactly zero.
  \item The mediation result: the effect of error magnitude on sequential
    Bayesian optimization acts through a one-shot selection loss computable from
    the landscape alone, at a cost of a fraction of a second.
  \item Two ablations that rule one confound out and one in. Telling the
    surrogate its true noise level removes only $5\%$ of the measured cost and
    nothing survives correction, so the cost is an information cost; changing
    the incumbent definition removes $15.6\%$ overall and $35$--$41\%$ for the
    probability-of-improvement family, so acquisition robustness rankings are
    really rankings of acquisition-and-incumbent pairs.
  \item Purpose-built landscape families that vary signal strength, sparsity and
    dimension independently, showing that the observational descriptor
    regression over-reads one predictor and that a within-family ladder indexed
    by dimension can still be measuring signal strength.
  \item A head-to-head run of the fitted-oracle design on the same grid. It
    reports roughly a third of the cost, and a manipulation worth up to $0.37$
    held-out $R^2$ in oracle accuracy barely moves it, so the discrepancy is
    attributable to the design rather than to surrogate fidelity.
  \item Generalisation checks that each change one thing: seven multi-objective
    problems with published hypervolume optima, budgets of $25$ and $100$, a
    bounded and discretised rating scale, two robust baselines, and ten
    preregistered held-out seeds.
\end{enumerate}

\section{Related work}

\paragraph{Bayesian optimization under observation noise.}
Noise has been part of the formulation since \citet{jones1998efficient} and is
routinely present in the settings that popularised the method
\citep{snoek2012practical}, and
several acquisitions are built for it: augmented expected improvement
\citep{huang2006global}, the knowledge gradient \citep{frazier2009knowledge},
and noisy expected improvement \citep{letham2019constrained}, whose Monte-Carlo
form is standard in modern implementations \citep{balandat2020botorch}.
\citet{picheny2013benchmark} benchmark infill criteria under noise on analytic
problems, which is the closest antecedent to what we do here; we differ in
scale, in treating the error \emph{process} rather than only its variance as an
experimental factor, and in asking which landscape properties predict the
outcome rather than only which criterion wins. Heteroscedastic surrogates
\citep{kersting2007most} address a case we deliberately do not cover: our error
is homoscedastic within a run, but arrives partway through it.

\paragraph{Human-in-the-loop and preference-based optimization.}
Where the objective is a person's judgement, the usual response is to change the
observation model rather than the noise model: elicit pairwise preferences and
fit an ordinal likelihood \citep{chu2005preference, gonzalez2017preferential}.
That is a good answer to inconsistent \emph{scaling} and a poor one to
inconsistent \emph{ranking}, and it does not tell a practitioner running a
cardinal-rating study what their measured rater noise will cost. The applied
literature that does ask directly builds the fitted-oracle simulation described
above. We take the complementary position: give up realism of the objective
entirely in exchange for exactness, and recover interpretability at the end by
mapping measured rater noise onto the synthetic axis.

\paragraph{Benchmarking practice.}
Analytic test functions with known optima are the field's default for algorithm
comparison, and the accompanying critique --- that results on them need not
transfer, and that suites are small and idiosyncratic --- is well made
\citep{eggensperger2021hpobench}. Our use of them is different in kind. We are
not ranking optimizers on a benchmark; we are using twenty landscapes as twenty
\emph{observations} in a regression whose unit of analysis is the landscape, and
then manipulating landscape geometry directly when the observational design
cannot identify an effect. Under that use, the heterogeneity of a hand-assembled
suite is the design's strength rather than its weakness.

\paragraph{Fitness landscape analysis.}
Characterising a problem by cheap statistics and predicting algorithm behaviour
from them is an old programme in evolutionary computation: autocorrelation and
ruggedness \citep{weinberger1990correlated}, fitness-distance correlation
\citep{jones1995fitness}, and the many descendants surveyed by
\citet{malan2013survey}. Our mediator $\frag(\sigma_e)$ belongs to that family
but differs in being defined \emph{relative to the perturbation of interest}
rather than to the landscape alone: it is not a property of the problem but of
the problem-and-error-level pair, which is what makes it predictive across both
axes at once.

\section{Setup}

\subsection{Landscapes}
We use twenty analytic benchmarks spanning $d = 2$ to $d = 11$: eight classical
functions (Schwefel, Powell, Eggholder, Ackley, Shekel, Griewank, Hartmann-3,
Hartmann-6), four further standard problems (Branin, Rosenbrock, Rastrigin,
Michalewicz), five psychophysical response functions including the
Yerkes--Dodson inverted-U \citep{yerkes1908relation}, and three
dynamic-optimization landscapes including the moving-peaks benchmark
\citep{branke1999memory}. Table~\ref{tab:benchmarks} reports their measured
geometry. The suite was assembled for a separate study and adopted unchanged, so
its composition is not tuned to the questions asked here.

Each objective is affinely standardised to zero mean and unit variance over its
own box using a fixed scrambled-Sobol sample of $2^{16}$ points. Raw outputs
span eight orders of magnitude across the suite, so without this a single grid
of error magnitudes would be imperceptible on one landscape and overwhelming on
the next. The map is strictly increasing, so the optimization problem itself is
unchanged: same argmax, same ordering of designs.

One scale is deliberately left \emph{un}normalised. Write $\optz$ for the number
of landscape standard deviations by which the optimum exceeds the mean --- that
is, the expected regret of a design drawn at random. Across the suite $\optz$
ranges from $1.25$ to $56.8$, so the same nominal error is trivial on one
landscape and catastrophic on another. That $45\times$ spread is what identifies
which scale the cost of error actually follows (Section~\ref{sec:currency});
normalising it away would remove the only variation that can answer the question.

\begin{table}[htbp]
\centering
\small
\caption{The twenty landscapes and their measured geometry. $\mathrm{opt}_z$ is
how many landscape standard deviations the optimum stands above a random design;
sparsity is the fraction of the box within 10\% of the optimum; ruggedness is
$1$ minus the lag-1 autocorrelation along random walks; tail is
$\sigma/(1.4826\,\mathrm{MAD})$; $\mathrm{frag}(1)$ is the one-shot selection
loss at an error of one landscape standard deviation. Names are the suite's
keys; \texttt{levy\_10} has eleven inputs.}
\label{tab:benchmarks}
\begin{tabular}{lrrrrrrr}
\toprule
landscape & $d$ & $\mathrm{opt}_z$ & sparsity & rugged. & skew & tail & $\mathrm{frag}(1)$ \\
\midrule
\texttt{powell} & 4 & 0.76 & 7.2e-01 & 0.13 & $-$2.33 & 2.52 & 0.18 \\
\texttt{rosenbrock} & 4 & 1.03 & 4.9e-01 & 0.13 & $-$1.06 & 1.13 & 0.18 \\
\texttt{branin} & 2 & 1.25 & 2.4e-01 & 0.13 & $-$1.09 & 1.14 & 0.30 \\
\texttt{steering\_law} & 4 & 1.29 & 7.9e-01 & 0.12 & $-$3.33 & 1.44 & 0.43 \\
\texttt{power\_law\_practice} & 4 & 1.72 & 1.4e-01 & 0.13 & $-$2.02 & 1.72 & 0.63 \\
\texttt{weber\_fechner} & 4 & 2.15 & 5.5e-02 & 0.15 & $-$0.71 & 1.01 & 0.64 \\
\texttt{levy\_10} & 11 & 2.62 & 1.3e-02 & 0.20 & $-$0.54 & 1.00 & 0.80 \\
\texttt{griewank} & 6 & 2.75 & 1.0e-02 & 0.12 & $-$0.26 & 0.98 & 0.85 \\
\texttt{moving\_peaks} & 4 & 2.96 & 1.8e-03 & 0.17 & $-$0.53 & 1.04 & 0.64 \\
\texttt{hartmann\_3} & 3 & 3.05 & 1.6e-02 & 0.14 & +1.14 & 1.49 & 0.50 \\
\texttt{stevens} & 4 & 3.06 & 3.0e-03 & 0.12 & $-$0.17 & 0.97 & 0.75 \\
\texttt{eggholder} & 2 & 3.21 & 6.6e-03 & 0.42 & +0.04 & 1.00 & 0.68 \\
\texttt{yerkes\_dodson} & 4 & 3.54 & 1.9e-03 & 0.12 & +0.27 & 0.98 & 0.73 \\
\texttt{rastrigin} & 4 & 3.64 & 5.8e-04 & 0.58 & $-$0.08 & 0.98 & 0.69 \\
\texttt{schwefel} & 4 & 4.33 & 1.2e-04 & 0.25 & +0.00 & 0.99 & 0.69 \\
\texttt{hicks\_law} & 4 & 5.40 & $<$1.5e-05 & 0.12 & +0.40 & 0.97 & 1.00 \\
\texttt{hartmann\_6} & 6 & 7.97 & 9.2e-05 & 0.14 & +2.67 & 2.74 & 0.53 \\
\texttt{michalewicz} & 5 & 8.05 & $<$1.5e-05 & 0.36 & +1.02 & 0.92 & 0.80 \\
\texttt{ackley} & 4 & 20.10 & $<$1.5e-05 & 0.80 & +3.26 & 1.67 & 0.18 \\
\texttt{shekel} & 4 & 56.77 & $<$1.5e-05 & 0.13 & +3.72 & 1.63 & 0.00 \\
\bottomrule
\end{tabular}

\end{table}


\subsection{Error processes}
\label{sec:errors}
Let $g(x)$ be the standardised objective, $t$ the iteration, $t_0$ the onset and
$\sigma_e$ the swept magnitude, in landscape standard deviations. Observations
are exact for $t \le t_0$; thereafter the optimizer sees $\tilde g_t$:
\begin{align}
  \text{\textsc{gaussian}}\quad & \tilde g_t = g(x_t) + \epsilon_t,
    & \epsilon_t &\sim \mathcal{N}(0, \sigma_e^2) \\
  \text{\textsc{bias}}\quad & \tilde g_t = g(x_t) + \sigma_e + \epsilon_t,
    & \epsilon_t &\sim \mathcal{N}(0, \sigma_e^2) \\
  \text{\textsc{drift}}\quad & \tilde g_t = g(x_t)
    + \sigma_e\tfrac{t - t_0}{T - t_0} + \epsilon_t,
    & \epsilon_t &\sim \mathcal{N}(0, \sigma_e^2) \\
  \text{\textsc{ar}(1)}\quad & \tilde g_t = g(x_t) + e_t,
    \quad e_t = \rho e_{t-1} + \eta_t,
    & \eta_t &\sim \mathcal{N}\!\left(0, \sigma_e^2(1 - \rho^2)\right)
\end{align}
with $\rho = 0.8$ and $e_{t_0+1}$ drawn from the stationary distribution, so the
AR(1) process has standard deviation $\sigma_e$ throughout rather than starting
cold. \textsc{bias} models a systematically generous or harsh rater, and its
offset is scaled to the swept magnitude so that it and \textsc{gaussian} contrast
a systematic against a random error \emph{of the same size}; with the offset held
at a fixed constant while the random component sweeps, as is otherwise natural,
the two conditions become indistinguishable at the top of the grid and the
contrast is untestable. \textsc{drift} models fatigue accumulating over a
session; \textsc{ar}(1) models a rater whose consecutive judgements are
correlated.

Noise first affects the observation at $t_0 + 1$; the first candidate that can
react to it is proposed at $t_0 + 2$. We sweep
$\sigma_e \in \{0.05, 0.25, 1, 5\}$ and $t_0 \in \{0, 20\}$ over $T = 50$
iterations with $5$ initial quasi-random samples, twelve acquisition functions
and ten seeds: $79{,}200$ runs, of which $76{,}800$ form noisy/clean pairs.

\subsection{Optimizer and acquisitions}
Every arm uses the same surrogate: a single-task GP with a Matern kernel,
inputs normalised to the unit cube and outputs standardised, fitted by marginal
likelihood, as implemented in BoTorch \citep{balandat2020botorch}. Only the
acquisition function varies. We run ten model-based arms --- expected
improvement and probability of improvement in their analytic, log
\citep{ament2023unexpected} and Monte-Carlo forms (EI, LogEI, qEI, PI, LogPI,
qPI), the upper confidence bound in analytic and Monte-Carlo form
\citep{srinivas2010gaussian} (UCB, qUCB), noisy expected improvement
\citep{letham2019constrained} (qNEI), and pure exploitation of the posterior
mean (Greedy) --- together with two model-free floors, uniform random sampling
and a scrambled Sobol sequence. Improvement-based acquisitions take the maximum
posterior mean at visited points as their incumbent, which is the noise-robust
choice and, as Appendix~\ref{sec:incumbent} shows, not a neutral one.

\subsection{Metrics}
The primary response is post-onset per-iteration excess simple regret in
landscape standard deviations: the regret of the noisy run minus that of the
identically-seeded clean run, averaged over the iterations the error could have
affected. Baseline and noisy runs share their candidate-proposal random stream,
so the pairing isolates the effect of the error rather than of the seed. Because
$y^\star$ is a verified supremum rather than a random-search estimate, this is a
genuine regret and cannot be negative by construction.

For reporting we also use the same quantity divided by $\optz$ --- the fraction
of the achievable improvement the error destroys, which is the number a
practitioner cares about. We never regress that normalised form on $\optz$, or
on anything collinear with it. Dividing by $\optz$ and then regressing on
$\log\optz$ manufactures the coefficient Section~\ref{sec:currency} exists to
estimate: if the raw excess were perfectly landscape-invariant, the normalised
response alone would still return the answer ``the achievable gain is the
currency'', by arithmetic rather than by evidence.

\subsection{Controls}
\label{sec:controls}
\begin{itemize}
  \item \textbf{Model-free floor.} \texttt{Random} and \texttt{Sobol} select
    candidates without reading any observation, so their excess regret must be
    identically zero. Measured across all $12{,}800$ such comparisons: exactly
    $0.000\mathrm{e}{+}00$. This is a positive control on the whole pipeline ---
    it establishes that no observation noise leaks into the candidate stream by
    any route --- and the two arms are then excluded from the robustness
    ranking, since a method that ignores its data wins such a ranking outright
    while learning nothing.
  \item \textbf{Headroom.} A landscape on which the clean run already reaches
    the optimum cannot show an effect, so a null on it is guaranteed by
    construction rather than earned. The lowest clean-run advantage over the
    model-free floor is $0.176$ (Rosenbrock); no landscape is ceiling-limited.
  \item \textbf{Uniform offset.} A constant added to every observation is
    algebraically invisible to a GP that standardises its targets and uses a
    posterior-mean incumbent. Driven with a shared noise stream this holds to
    $9 \times 10^{-12}$ over the initial design and the first three model-based
    candidates, after which floating-point differences in the standardisation
    are amplified chaotically by the acquisition optimizer. In the sweep the
    \textsc{bias}-at-onset-0 and \textsc{gaussian}-at-onset-0 arms agree to
    within $0.011$ at every magnitude, as they should.
\end{itemize}

\section{How much does feedback error cost?}
\label{sec:cost}

\begin{table}[t]
\centering
\caption{Fraction of the achievable improvement destroyed by feedback error,
averaged over landscapes and model-based acquisitions. Magnitudes are in
landscape standard deviations.}
\label{tab:dose}
\begin{tabular}{lrrrr}
\toprule
error present & $0.05\sigma$ & $0.25\sigma$ & $1\sigma$ & $5\sigma$ \\
\midrule
from it.\ 1 & 0.9\% & 5.3\% & 14.1\% & 22.0\% \\
from it.\ 21 & 0.3\% & 1.2\% & 2.8\% & 4.9\% \\
\bottomrule
\end{tabular}

\end{table}

Table~\ref{tab:dose} is the headline. Error whose standard deviation equals the
landscape's own, present from the first rating, costs $14.1\%$ of the
improvement that was available; even error five times that large costs $22.0\%$.
This is considerably less damage than the framing of the applied literature
anticipates, and the reason is visible in the second row: an optimizer that has
already located a good region is very hard to dislodge. The same $1\sigma$ error
arriving at iteration 21 rather than iteration 1 costs $2.8\%$ instead of
$14.1\%$, a factor of five.

Decomposing the variation by factor puts the practical conclusions in order. At
a fixed onset, sweeping the magnitude over its two-decade range moves the cost
by $23\times$; changing the landscape at fixed magnitude moves it by $8.1\times$
(Powell $0.038$ to Hartmann-6 $0.303$ at $1\sigma$); moving the onset moves it by
$5.0\times$; changing the acquisition function moves it by $2.8\times$; and
changing the error \emph{process} moves it by $1.32\times$. The two factors a
practitioner controls --- which acquisition to run, and which error model to
worry about --- are the two smallest, and the two that dominate are properties of
the problem and of the raters.

That the error process barely matters deserves emphasis, because a good deal of
methodological effort goes into distinguishing kinds of human error. At matched
magnitude and onset~0 the mean cost is $0.109$ for gaussian, $0.113$ for drift,
$0.110$ for systematic bias and $0.092$ for AR(1). Systematic error is
indistinguishable from random error of the same size, and serially correlated
error is if anything \emph{gentler} than white noise --- a correlated error
sequence is partly absorbed into the surrogate's mean rather than scattering
its observations. For the purposes of predicting damage, the magnitude and the
timing are the error's only relevant properties.

\paragraph{Is $1\sigma$ a lot?} A landscape standard deviation is only a useful
unit if real feedback noise can be expressed in it. We estimate within-rater
noise in three published human studies with a nearest-neighbour (semivariogram
nugget) estimator over repeated ratings by the same participant, and divide it
by the standard deviation of that study's own fitted objective over its design
box --- the same quantity $\sigma_f$ denotes for a benchmark
(Table~\ref{tab:anchor}).

Real rating noise lands between $0.74$ and $3.35$ landscape standard deviations.
Participants are noisy from their first rating, so those studies sit in the
onset-0 row of Table~\ref{tab:dose} between its $1\sigma$ and $5\sigma$ columns:
a $12$--$20\%$ loss of the achievable improvement. The magnitudes at which we
measure the largest costs are the magnitudes human raters actually produce, so
the robustness reported here is a genuine finding about the applied setting
rather than an artefact of a benign choice of grid. The caveat is structural:
$\sigma_f$ for a real study can only be computed through a fitted oracle, which
is precisely the object the synthetic arm exists to avoid. This anchors the
interesting range; it is not itself a measurement.

\begin{table}[t]
\centering
\caption{Three published human studies placed on the synthetic arm's axis.
$\sigma_f$ is the spread of the study's fitted objective over its design box;
the last column is the measured within-rater noise in those units.}
\label{tab:anchor}
\begin{tabular}{llrrr}
\toprule
study & ratings & $\sigma_f$ & noise / $\sigma_f$ & upper bound \\
\midrule
\texttt{ehmi} & as logged & 0.365 & 0.74 {\scriptsize $[0.48, 0.97]$} & 1.17 \\
\texttt{opticarvis} & as logged, two scales mixed & 0.374 & 3.35 {\scriptsize $[3.03, 3.58]$} & -- \\
\texttt{opticarvis} & one scale & 0.156 & 1.11 {\scriptsize $[0.81, 1.40]$} & 1.34 \\
\texttt{provoice} & as logged, one sign reversed & 0.298 & 1.02 {\scriptsize $[0.20, 1.56]$} (weak) & 2.58 \\
\texttt{provoice} & sign corrected & 0.449 & 0.60 {\scriptsize $[0.00, 0.90]$} (weak) & 2.37 \\
\bottomrule
\end{tabular}

\end{table}

\section{What makes an error big?}
\label{sec:currency}

Standardising every objective makes the sweep possible, but it embeds an
assumption: that a landscape's own standard deviation is the right unit for its
error. There is an obvious rival. What might matter is the error relative to the
improvement actually available, $\sigma_e/\optz$ --- a half-SD error is one thing
when the optimum stands $1.25$ SDs above a random design and quite another when
it stands $56.8$ above. Regressing raw excess regret $R$ on both quantities,
\begin{equation}
  R \;=\; \beta_c \log_{10}\sigma_e \;+\; \beta_z \log_{10}\optz \;+\; \varepsilon,
\end{equation}
\textsc{spread} predicts $\beta_z = 0$ and \textsc{gain} predicts
$\beta_z = 1 - \beta_c$. Inference is a bootstrap that resamples \emph{landscapes},
which is the unit of replication for any landscape-level covariate; with twenty
clusters the asymptotic cluster-robust sandwich is anti-conservative and can come
out rank-deficient.

\begin{table}[t]
\centering
\caption{Neither scale is right at every onset. $\beta_z$ excludes zero in all
eight conditions, so standardising by the landscape's spread is not sufficient;
at a mid-run onset the achievable-gain scaling fits, but from the first
observation the cost grows faster than either prediction.}
\label{tab:currency}
\begin{tabular}{llrrcrcl}
\toprule
& & & \multicolumn{2}{c}{\textsc{spread}: $\beta_z = 0$} & \multicolumn{2}{c}{\textsc{gain}: $\beta_c + \beta_z = 1$} & \\
\cmidrule(lr){4-5} \cmidrule(lr){6-7}
error process & $t_0$ & $\beta_c$ & $\beta_z$ & 95\% CI & $\beta_c + \beta_z - 1$ & 95\% CI & verdict \\
\midrule
\textsc{gaussian} & 0 & 0.51 & 0.86 & [0.47, 1.13] & +0.37 & [+0.15, +0.70] & neither \\
\textsc{bias} & 0 & 0.51 & 0.85 & [0.49, 1.08] & +0.36 & [+0.17, +0.65] & neither \\
\textsc{drift} & 0 & 0.55 & 0.82 & [0.55, 1.11] & +0.37 & [$-$0.01, +0.63] & gain \\
\textsc{ar}(1) & 0 & 0.66 & 0.85 & [0.60, 1.15] & +0.51 & [+0.14, +0.77] & neither \\
\textsc{gaussian} & 20 & 0.62 & 1.15 & [0.64, 1.48] & +0.76 & [+0.03, +1.03] & neither \\
\textsc{bias} & 20 & 0.70 & 1.14 & [0.61, 1.58] & +0.85 & [+0.04, +1.08] & neither \\
\textsc{drift} & 20 & 0.65 & 1.15 & [0.61, 1.55] & +0.80 & [$-$0.00, +1.14] & gain \\
\textsc{ar}(1) & 20 & 0.71 & 1.14 & [0.64, 1.49] & +0.85 & [+0.13, +1.11] & neither \\
\bottomrule
\end{tabular}

\end{table}

Table~\ref{tab:currency} rejects \textsc{spread} outright: $\beta_z$ is positive
with an interval excluding zero in all eight conditions, so standardising an
objective is \emph{not} enough to make error magnitudes comparable across
problems, and any study that reports a noise level without reporting how much
headroom the problem had is under-specified. At a mid-run onset the
achievable-gain reading fits: $\beta_z$ lands within its interval of $1 -
\beta_c$ for all four error processes.

From the first observation it does not. There $\beta_z$ runs $1.70$ to $2.05$
against a \textsc{gain} prediction of $0.26$ to $0.30$: the cost grows
\emph{superlinearly} in how far the optimum stands above a random design. The
natural reading is that these are two different failures. A late-arriving error
degrades an optimizer that already knows roughly where to look, and costs it a
share of the remaining gap --- the \textsc{gain} scaling. An error present from
the start degrades the \emph{search}, and a tall isolated optimum is
disproportionately expensive to locate through noise, because the evidence
distinguishing it from its surroundings is exactly what the noise destroys.
Appendix~\ref{sec:manipulation} tests this reading directly and qualifies it.

\section{The cost is mediated by a one-shot selection loss}
\label{sec:mediation}

Both preceding sections describe the cost of error in terms of properties of the
landscape, but neither gives a practitioner a number. This section does. Define
$\frag(\sigma_e)$ as the expected loss from a single greedy choice among $m$
quasi-random candidates when each is observed with $\mathcal{N}(0,\sigma_e^2)$
error:
\begin{equation}
  \frag(\sigma_e) \;=\;
  \mathbb{E}_{\epsilon}\Big[\max_i g_i - g_{\arg\max_i (g_i + \epsilon_i)}\Big],
  \qquad \epsilon_i \sim \mathcal{N}(0,\sigma_e^2),
  \label{eq:frag}
\end{equation}
where $g_i = g(x_i)$ over a fixed Sobol design. It involves no surrogate, no
sequence and no acquisition function; at $m = 256$ and $4000$ Monte-Carlo draws
it costs about $0.2$ seconds per landscape. Being indexed by the error level as
well as the landscape, it varies along both axes of the sweep.

\begin{table}[t]
\centering
\caption{Mediation. Fitted on landscape $\times$ condition cell means,
predictors $z$-scored, intervals and $p$-values from a 2000-replicate bootstrap
over landscapes. Conditional on the one-shot selection loss, the error magnitude
explains nothing.}
\label{tab:mediation}
% Seven numeric columns with long math headers run 22pt past the text width at
% normal size. Set here rather than in tables/mediation.tex, which is generated
% and would lose the fix on the next rebuild.
\small
\setlength{\tabcolsep}{5pt}
\begin{tabular}{lrrrrrrrr}
\toprule
model & $\mathrm{frag}(\sigma_e)$ & $\log\sigma_e$ & $\log\mathrm{opt}_z$ & $\log\sigma_e{\times}\log\mathrm{opt}_z$ & $\log$ tail & rugged. & $d$ & $R^2$ \\
\midrule
indicators only & -- & -- & -- & -- & -- & -- & -- & 0.040 \\
$\log\sigma_e$ only & -- & +0.37$^{***}$ & -- & -- & -- & -- & -- & 0.158 \\
$\mathrm{frag}$ only & +0.75$^{***}$ & -- & -- & -- & -- & -- & -- & 0.520 \\
descriptors only & -- & +0.37$^{***}$ & +0.64$^{***}$ & -- & +0.20$^{**}$ & $-$0.09 & $-$0.03 & 0.526 \\
both & +0.66$^{***}$ & $-$0.07 & +0.42$^{**}$ & -- & +0.20$^{**}$ & $-$0.06 & $-$0.03 & 0.701 \\
descriptors + int. & -- & +0.37$^{***}$ & +0.64$^{***}$ & +0.54$^{***}$ & +0.20$^{**}$ & $-$0.09 & $-$0.03 & 0.779 \\
both + int. & +0.31$^{***}$ & +0.17$^{*}$ & +0.54$^{**}$ & +0.40$^{***}$ & +0.20$^{**}$ & $-$0.08 & $-$0.03 & 0.792 \\
\bottomrule
\end{tabular}

\end{table}

Table~\ref{tab:mediation} is the paper's central result. On its own,
$\frag(\sigma_e)$ explains $R^2 = 0.52$ of the variation in excess regret across
landscape-by-condition cells --- essentially as much as the entire set of
pre-measured landscape descriptors together with the error magnitude
($R^2 = 0.53$). Entered jointly the two are complementary ($R^2 = 0.70$), and
the informative part is which coefficient survives: $\frag$ keeps
$\beta = +0.66$ with a bootstrap interval of $[0.28, 0.74]$, while the error
magnitude collapses to $\beta = -0.07$ with $p = 0.44$. \textbf{Conditional on
what one noisy choice would cost on this landscape, how large the error is
carries no further information about what fifty sequential noisy choices will
cost.} Dimension and ruggedness are null throughout, in this analysis, though
Appendix~\ref{sec:manipulation}
shows that the null on dimension is an artefact of the observational design.

The practical consequence is the reason to care. Equation~\ref{eq:frag} needs
only a design space and an estimate of rater noise, both of which a practitioner
planning a study has before they run it. The residual terms say where the
one-shot approximation breaks down: $\optz$ and tail weight both retain
significant coefficients, which is where a surrogate fitted to fifty
observations averages noise down in a way a single choice cannot.

\section{Which acquisition function?}
\label{sec:acq}

\begin{table}[t]
\centering
\caption{Acquisition robustness, ranked over landscapes within each condition
and averaged over conditions. ``Absolute loss'' is the deployment number: the
post-onset regret under noise as a fraction of the achievable gain.}
\label{tab:acq}
\begin{tabular}{lrrr}
\toprule
acquisition & mean rank & excess regret & absolute loss \\
\midrule
qUCB & 3.95 & 4.14\% & 27.9\% \\
qNEI & 4.34 & 4.89\% & 25.2\% \\
UCB & 4.41 & 4.73\% & 27.4\% \\
LogEI & 4.61 & 5.21\% & 26.0\% \\
qEI & 4.73 & 5.37\% & 25.5\% \\
EI & 4.98 & 5.42\% & 26.2\% \\
qPI & 6.29 & 7.90\% & 31.0\% \\
Greedy & 6.56 & 8.11\% & 32.5\% \\
LogPI & 7.48 & 9.19\% & 31.3\% \\
PI & 7.66 & 9.44\% & 31.7\% \\
\midrule
model-free floor & -- & 0.00\% & 41.0\% \\
\bottomrule
\end{tabular}

\end{table}

The confidence-bound family and noisy expected improvement are the most robust,
and the probability-of-improvement family is uniformly the worst, on both
robustness and absolute performance (Table~\ref{tab:acq}). The ordering is not
an artefact of the robustness metric: the same acquisitions lead on absolute
loss under noise, and every model-based arm is far better than the model-free
floor at $0.410$, so no arm is winning by declining to learn. The Friedman test
over the twenty landscapes as blocks reaches $p < 0.05$ in $28$ of $32$
conditions, and $115$ of $288$ pairwise comparisons against the per-condition
winner survive Benjamini--Hochberg correction \citep{benjamini1995controlling}.

The result that matters more is the disagreement. Kendall's $W$ across
landscapes is $0.247$: the twenty landscapes largely do not agree on an ordering.
qUCB takes $10$ of $32$ conditions, UCB and qNEI $6$ each, LogEI $5$, EI $3$,
qEI $2$. A study that measured any one of these landscapes and reported its
winner would have produced a defensible, statistically significant, and largely
non-transferable recommendation. Running on one design space, the applied
literature is structurally obliged to do exactly that. Appendix~\ref{sec:incumbent}
shows that a further part of this ranking is not about the acquisitions at all.

\section{What the fitted-oracle design would have concluded}
\label{sec:fitted}

The introduction argued that fitting a regression oracle to archival ratings
entangles the cost of feedback error with the cost of the oracle being wrong.
That is an argument, not a measurement, until the two designs are run on a
matched grid. We run the data-driven simulator on the three
archival datasets at the same error magnitudes, expressed in each dataset's own
$\sigma_f$ so that ``$1\sigma$'' means the same thing in both arms, and compare.

\begin{table}[t]
\centering
\caption{The same experiment, run against exact functions and against oracles
fitted to ratings. Percentages are post-onset excess as a fraction of the gap
between the optimum and a same-budget model-free floor --- the one normaliser
computable in both arms. The fitted design reports roughly a third of the cost.}
\label{tab:fitted}
\begin{tabular}{lrrrrrrrr}
\toprule
& \multicolumn{4}{c}{error from it.\ 1} & \multicolumn{4}{c}{error from it.\ 21} \\
\cmidrule(lr){2-5} \cmidrule(lr){6-9}
arm & $0.05\sigma$ & $0.25\sigma$ & $1\sigma$ & $5\sigma$ & $0.05\sigma$ & $0.25\sigma$ & $1\sigma$ & $5\sigma$ \\
\midrule
known functions (20) & 10 & 44 & 99 & 148 & 3 & 7 & 10 & 16 \\
fitted oracle (3) & 1 & 11 & 33 & 47 & 1 & 3 & 6 & 8 \\
fitted, no augmentation (3) & 2 & 14 & 38 & 50 & 1 & 4 & 6 & 7 \\
\bottomrule
\end{tabular}

\end{table}

Two things make the comparison awkward. The fitted arm has no verified optimum,
so its reference is the best value any arm reached --- attainable by
construction, which shrinks its denominator and inflates its reported fraction
relative to what a true supremum would give. That one is conservative: it works
against the result below. The second is not conservative but is a limit on
precision rather than a bias --- three design spaces against twenty, so the
fitted intervals are wide.

\textbf{The fitted design understates the cost of feedback error by roughly
threefold.} At $1\sigma$ from the first rating the known-function arm puts the
loss at $99\%$ of the achievable gain, with a cluster-bootstrap interval of
$[55, 155]$; the fitted arm puts it at $33\%$, interval $[5, 75]$. At
$0.05\sigma$ the fitted arm reports $1\%$, interval $[-4, 3]$ --- indistinguishable
from no effect --- where the exact functions report $10\%$, interval $[4, 18]$.
The direction is the same at all four magnitudes and both onsets: a factor of
three at $1\sigma$ and $5\sigma$, four at $0.25\sigma$, and at $0.05\sigma$ a
ratio we will not quote, because the fitted arm's estimate there is not
distinguishable from zero and the quotient is unstable. With three datasets the
intervals overlap at the larger magnitudes, so this is a consistent gap rather
than a separation at every cell.

\textbf{Oracle fidelity does not explain it.} The obvious reading is that the
fitted oracles are simply poor, and the deployed jitter augmentation gives a
handle on that: turning it off is worth up to $0.37$ held-out $R^2$. We reran the
whole companion arm with augmentation disabled. Paired on the $2{,}400$ cells the
two arms share, the change is $+0.016$ in the fraction destroyed (median
$+0.006$), and the augmentation-free arm is the worse of the two in $52.6\%$ of
cells --- a coin flip. The oracles top out near $R^2 = 0.55$, so a manipulation
worth up to $0.37$ moves most of what they explain, and it moves the measured
cost of feedback error almost not at all.

The negative is the useful half. If the gap were fidelity it would close as the
oracle improved, and the fix for the applied literature would be better oracles.
It does not behave that way here. What the fitted design changes is not how accurately the
surrogate human answers, but what the regret is measured \emph{against}: a
random-search estimate of the optimum that the optimizer reaches routinely
compresses the range in which any error can be seen to do damage. The onset
effect compresses with it --- $0.8$ to $5.6\times$ across magnitudes in the fitted
arm against $3.6$ to $9.6\times$ in the exact one.

\section{When the error is in the design, not the rating}
\label{sec:inputerror}

Every error process above corrupts the reported value: the optimizer is told the
wrong number about the design it proposed. People also make the other mistake.
A slider overshoots, the wrong option is pressed, a configuration is applied to
the wrong condition. The optimizer proposes $x$, the person experiences
$x' \neq x$, rates it honestly, and the rating is filed against $x$. A response
error is a wrong value at the right location; this is the right value at the
wrong location.

It is also the corruption a fitted-oracle design is least able to measure.
Scoring it requires the objective at the point actually touched. A fitted oracle
supplies $\hat g(x')$, so the measured damage is the slip's cost plus the
oracle's error at a second location. On an exact landscape, $g(x')$ is exact.

We model two input processes. A \textsc{slip} displaces every trial by
$\delta \sim \mathcal{N}(0, s^2)$ per coordinate, with $s$ a fraction of that
coordinate's range, and clips the result to the box, as a bounded control stops
at its end. A \textsc{misclick} leaves the trial exact with probability $1 - p$
and otherwise lands on a uniformly random design. Both magnitudes are
dimensionless and swept over $\{0.01, 0.05, 0.15, 0.4\}$; neither is in
landscape standard deviations, so the numbers in this section are not comparable
with the response-error grid. The rating itself is exact throughout.

What gets written down matters as much as where the person went. In the
realistic case the slip goes unnoticed, the surrogate trains on $(x, g(x'))$,
and the design a practitioner eventually deploys is the recorded $x$ --- never
the $x'$ that earned its rating. We therefore score the deployed design at the
recorded point. A counterfactual arm logs $(x', g(x'))$ instead, as an
instrumented interface would, with identical slip draws; the contrast separates
the cost of going to the wrong place from the cost of mislabelling it.
Model-free floors run alongside, and here they are not a zero control: they
never read an observation, but they do evaluate the wrong points, so their
excess regret is the pure geometric cost of a slip with no learning to corrupt.

\todo{Results, once the three input-error arms complete: dose-response by slip
and misclick magnitude; the recorded-versus-actual contrast; model-based excess
against the floors' geometric cost; deployed-design regret against evaluated
regret.}

\section{The conclusions under six further checks}
\label{sec:robustness}

Each check below changes one thing about the design; all six are reported in
full in the appendices.

\paragraph{The surrogate's noise estimate is not the problem.} Telling the GP the
true injected variance removes $5.0\%$ of the measured cost and no condition
survives correction, even though its fitted noise is an eighth of the true
value for the first twenty iterations. The cost is an information cost, and the
onset effect is not a failure of noise estimation
(Appendix~\ref{sec:noisediag}).

\paragraph{The incumbent is part of the method.} Switching improvement-based
acquisitions from the posterior-mean incumbent to the observed maximum removes
$15.6\%$ of the cost overall and $35$--$41\%$ for the probability-of-improvement
family, so a robustness ranking over acquisitions is a ranking over
acquisition-and-incumbent pairs (Appendix~\ref{sec:incumbent}).

\paragraph{Most of the apparent dimension effect is signal strength.} Landscapes
built to vary one property at a time show that $\optz$ is not causal on its own,
and holding $\optz$ fixed removes most of the dimension effect a Levy ladder
appeared to show (Appendix~\ref{sec:manipulation}).

\paragraph{Multi-objective BO is no more fragile.} On six \textsc{BoTorch} problems
with published hypervolume optima the dose-response and onset structures
replicate; on a common normaliser early error costs roughly half as much as on
the scalar suite, with overlapping intervals at the largest magnitudes, and the
four hypervolume acquisitions are indistinguishable
(Appendix~\ref{sec:multiobjective}).

\paragraph{The onset effect depends on the budget.} With the late onset held at
$40\%$ of the run, the early-to-late cost ratio is $2.3\times$, $5.0\times$ and
$11.7\times$ at $T = 25$, $50$ and $100$, so every onset number above describes a
fifty-iteration study. A bounded, discretised rating scale makes small errors
costlier, since rounding is itself an error, and large ones slightly cheaper;
qKG and replication, two methods built for noise, recover about a third of the
early-error cost (Appendix~\ref{sec:designchoices}).

\paragraph{The claims hold on held-out seeds, but the replication is void.} We
preregistered four claims before running ten fresh seeds. One control failed ---
a landscape fell below the preregistered headroom threshold --- so the run is void
by its own rule; with that landscape excluded post hoc, all four claims clear
their thresholds (Appendix~\ref{sec:confirmatory}).

\section{Discussion}
\label{sec:discussion}

The practical summary is short. Feedback error costs a real study something like
a tenth to a fifth of the improvement it could have had, which is worth reducing
but is not the catastrophe the framing sometimes implies. Effort spent
classifying \emph{what kind} of error one's raters make is largely wasted, since
the four processes we test are within $32\%$ of each other at matched magnitude.
Effort spent choosing an acquisition function is worth about $2.8\times$ in cost,
and is only transferable to the extent that a new design space resembles the one
it was measured on --- which, given $W = 0.247$, is not far. What does transfer
is Equation~\ref{eq:frag}: compute it for the design space and the measured rater
noise, and the expected damage follows.

The methodological finding may outlast the numbers. Running the fitted-oracle
design against the same error grid recovers the same qualitative structure at
about a third of the magnitude, and improving the oracle does not close the gap
(Section~\ref{sec:fitted}). A literature that measures robustness against a
random-search estimate of its own optimum is not measuring a smaller version of
the same thing; it is measuring a quantity whose ceiling it set itself. That is
worth knowing before the next such study is designed, and it is cheap to avoid.

Two things we could not explain are worth stating as open. The onset effect is
large and we do not know its mechanism: the obvious candidate, a surrogate that
cannot yet estimate its noise, was tested and rejected
(Appendix~\ref{sec:noisediag}), and the effect turns out to scale with the budget
(Appendix~\ref{sec:designchoices}) in a way a mechanism should explain. And the
superlinear growth of early-error cost in $\optz$ survives as a description
while the manipulation shows $\optz$ is not causally responsible on its own,
pointing at an interaction between signal strength and how findable the optimum
is that we have characterised but not modelled.

\section{Limitations}
\begin{itemize}
  \item The landscapes are analytic, not people. They remove the
    surrogate-human confound at the cost of the realism the fitted-oracle
    literature buys. Section~\ref{sec:fitted} runs both designs on a matched
    grid and finds the same qualitative structure at about a third of the
    magnitude, which establishes that the choice of design changes the answer
    but not which of the two answers is closer to a person.
  \item The anchoring of $\sigma_e$ to real rating noise runs through a fitted
    oracle's $\sigma_f$, so it inherits that oracle's error --- the very thing
    the synthetic design was built to avoid. It bounds the interesting range; it
    is not a measurement.
  \item The multi-objective arm (Appendix~\ref{sec:multiobjective}) is six
    problems at five seeds against the scalar arm's twenty landscapes at ten,
    so it can show that the dose-response and onset structures replicate but
    cannot separate the four hypervolume acquisitions, and its intervals at the
    top of the error range are wide. The mediator of
    Section~\ref{sec:mediation} has no Pareto analogue and is untested there.
  \item The robust baselines are two (Appendix~\ref{sec:designchoices}), not a
    survey. qKG and replication recover about a third of the early-error cost;
    whether a method built specifically for corrupted human feedback recovers
    more is untested.
  \item The onset ratio depends strongly on the budget --- $2.3\times$ at
    $T = 25$, $5.0\times$ at $T = 50$ and $11.7\times$ at $T = 100$
    (Appendix~\ref{sec:designchoices}). Every onset number in this paper is a
    statement about a fifty-iteration study.
  \item The preregistered replication is void on a control failure
    (Appendix~\ref{sec:confirmatory}). The four claims clear their thresholds on
    fresh seeds once the offending landscape is dropped, but that exclusion is
    post-hoc, so the strongest honest statement is that the claims survived an
    out-of-sample test committed to in advance and cannot be called confirmed.
\end{itemize}

\section*{Reproducibility statement}
The full pipeline is scripted end to end, from the landscape statistics through
the sweep to every table in this paper, and every stochastic component is
explicitly seeded. Regenerating published runs from their own recorded settings
reproduces every model-based run bit for bit. The one exception is the
model-free floors of the main sweep, whose noise seed moved when two robust
baselines were added to the acquisition list after that sweep ran: their
candidates and regret still reproduce exactly, since a floor never reads an
observation, but their recorded noisy observations do not. No reported number
depends on those, and the seed order is now pinned by test. Three specific
commitments are worth naming. The vendored
benchmark implementations are asserted by test to agree with their source
implementations to $10^{-9}$ relative on random points, so the objective cannot
silently drift. The model-free control of Section~\ref{sec:controls} is verified
to be exactly zero in every run, which is a stronger statement about pipeline
correctness than any test we could write. And one recorded optimum in the source
suite is stale --- it is the supremum for a superseded parameter setting, and is
wrong by $64\%$ --- which we detected by searching every box against its recorded
value and corrected here; the correction and the provenance of the original are
both asserted by test. \todo{repository URL and commit at camera-ready.}

\bibliographystyle{iclr2026_conference}
\bibliography{references}

\appendix
\section{Ablations}

\subsection{Does the surrogate know it is being lied to?}
\label{sec:noisediag}
The GP fits its observation noise as a free hyperparameter under a prior that
shrinks it, so a measured cost of noise could in principle be the surrogate's
failure to notice the error rather than the information the error destroys. This
is the first thing to rule out, because if it were the explanation then every
number above would be a statement about BoTorch's default prior rather than
about noisy feedback.

The mis-estimation is real, and it lands where it would do the most damage
(Table~\ref{tab:noisediag}). Refitting each surrogate from the sweep's own logs
using only its first $n$ observations, the fitted noise standard deviation is
about an eighth of the injected one for the first twenty of fifty iterations,
and has converged by the end of the run --- so a diagnostic evaluated only at
$n = 50$ looks reassuring and is misleading. Refitting with a wide noise prior
recovers the injected level far sooner ($0.55$ rather than $0.14$ at $n = 15$,
$1\sigma$), so the shrinkage is the prior's doing and not a limit of the data.

\begin{table}[t]
\centering
\caption{Fitted over injected observation-noise standard deviation, as a
function of how many observations the surrogate has seen, for error present from
iteration~1. The estimate is badly wrong exactly while the trajectory is being
decided --- and Table~\ref{tab:knownnoise} shows that this costs almost nothing.}
\label{tab:noisediag}
\begin{tabular}{lrrrr}
\toprule
observations & $0.05\sigma$ & $0.25\sigma$ & $1\sigma$ & $5\sigma$ \\
\midrule
8 & 1.43 & 0.30 & 0.12 & 0.11 \\
15 & 1.50 & 0.41 & 0.14 & 0.10 \\
25 & 1.37 & 0.72 & 0.66 & 0.56 \\
50 & 1.24 & 0.88 & 0.83 & 0.89 \\
\bottomrule
\end{tabular}

\end{table}

It was natural to expect this to explain the onset effect of
Section~\ref{sec:cost}: from the first observation the error arrives in exactly
the window where the noise cannot yet be estimated, whereas at a mid-run onset
the surrogate has a settled length scale and thirty iterations in which to
learn. \textbf{We tested that directly and it is wrong.} Supplying the true
injected variance as \texttt{train\_Yvar} and re-running $4{,}800$ paired cells
removes only $5.0\%$ of the measured cost, no condition survives FDR correction
(Table~\ref{tab:knownnoise}), and the onset-0 to onset-20 ratio at $1\sigma$ does
not shrink --- it goes from $4.0$ to $5.2$. Per landscape the largest reduction
is Ackley at $47\%$; per acquisition the effect is inconsistent in sign, helping
PI ($22\%$) and qNEI ($15\%$) but \emph{hurting} qUCB ($-29\%$).

Two conclusions follow. Bayesian optimization is remarkably insensitive to its
own noise hyperparameter in this regime, despite getting it badly wrong for the
first two-fifths of the budget --- worth knowing on its own, and a caution
against reading too much into noise-estimation diagnostics. And the cost this
paper measures is the information cost of the error, not an artefact of a
mis-specified surrogate, which is what licenses reading Table~\ref{tab:dose} as
a property of the problem. The onset effect itself is therefore left
unexplained; we return to it in Section~\ref{sec:discussion}.

\begin{table}[t]
\centering
\caption{Telling the surrogate the truth about its noise. Paired within
(landscape, acquisition, magnitude, onset, seed); positive $\Delta$ means the
treatment suffers more. Nothing survives FDR correction.}
\label{tab:knownnoise}
\begin{tabular}{rrrrrrrr}
\toprule
$\sigma_e$ & $t_0$ & reference & treatment & $\Delta$ & share & $d_z$ & $p_{\text{FDR}}$ \\
\midrule
0.05 & 0 & 0.065 & 0.077 & +0.012 & $-$19\% & 0.03 & 0.888 \\
0.25 & 0 & 0.271 & 0.166 & $-$0.105 & +39\% & $-$0.18 & 0.888 \\
1 & 0 & 0.790 & 0.848 & +0.058 & $-$7\% & 0.10 & 0.888 \\
5 & 0 & 1.484 & 1.463 & $-$0.021 & +1\% & $-$0.04 & 0.888 \\
0.05 & 20 & 0.002 & 0.001 & $-$0.002 & +71\% & $-$0.02 & 0.913 \\
0.25 & 20 & 0.072 & 0.041 & $-$0.031 & +43\% & $-$0.20 & 0.888 \\
1 & 20 & 0.197 & 0.164 & $-$0.034 & +17\% & $-$0.19 & 0.888 \\
5 & 20 & 0.531 & 0.483 & $-$0.048 & +9\% & $-$0.16 & 0.888 \\
\bottomrule
\end{tabular}

\end{table}

\subsection{Is it the acquisition or the incumbent?}
\label{sec:incumbent}
$\mathrm{best}_f$ is the maximum posterior mean at visited points, which shrinks
under noise and so makes improvement-based acquisitions automatically more
exploratory. The acquisitions differ in their exposure to this: LogEI, EI, PI,
LogPI and qEI consume $\mathrm{best}_f$; UCB does not; qNEI ignores it entirely.
A ranking of acquisitions by robustness is therefore partly a ranking of their
sensitivity to a design choice that is not part of the acquisition at all.
Re-running with the observed maximum as the incumbent removes $15.6\%$ of the
measured cost overall, with four of eight conditions significant after FDR
correction (Table~\ref{tab:incumbent}).

The per-acquisition breakdown (Table~\ref{tab:incumbentacq}) is the finding. UCB
and qNEI move by \emph{exactly} zero, as they must --- a control that validates
the pairing at machine precision. The probability-of-improvement family moves
most: LogPI loses $41\%$ of its measured cost and PI $35\%$, against $17\%$ for
EI, $2\%$ for LogEI and $-2\%$ for qEI. So a substantial part of ``PI is fragile
under noise'' is really ``the posterior-mean incumbent is fragile under noise,
and PI is the acquisition most exposed to it''. Every acquisition-robustness
ranking here, including Table~\ref{tab:acq}, is a ranking of
acquisition-and-incumbent pairs, and the same is true of the applied results it
is modelled on. Reporting the incumbent definition alongside the acquisition
should be standard.

\begin{table}[t]
\centering
\caption{Observed-maximum against posterior-mean incumbent, same pairing as
Table~\ref{tab:knownnoise}.}
\label{tab:incumbent}
\begin{tabular}{rrrrrrrr}
\toprule
$\sigma_e$ & $t_0$ & reference & treatment & $\Delta$ & share & $d_z$ & $p_{\text{FDR}}$ \\
\midrule
0.05 & 0 & 0.010 & 0.003 & $-$0.007 & +72\% & $-$0.56 & 0.019 \\
0.25 & 0 & 0.050 & 0.026 & $-$0.024 & +48\% & $-$1.39 & $<$0.001 \\
1 & 0 & 0.132 & 0.100 & $-$0.032 & +24\% & $-$0.93 & 0.004 \\
5 & 0 & 0.223 & 0.200 & $-$0.022 & +10\% & $-$0.71 & 0.011 \\
0.05 & 20 & 0.003 & 0.001 & $-$0.002 & +81\% & $-$0.24 & 0.142 \\
0.25 & 20 & 0.014 & 0.005 & $-$0.009 & +64\% & $-$0.88 & 0.002 \\
1 & 20 & 0.027 & 0.021 & $-$0.005 & +20\% & $-$0.50 & 0.032 \\
5 & 20 & 0.049 & 0.048 & $-$0.000 & +0\% & $-$0.01 & 0.421 \\
\bottomrule
\end{tabular}

\end{table}

\begin{table}[t]
\centering
\caption{The incumbent arm by acquisition. UCB and qNEI never read
$\mathrm{best}_f$ and move by exactly zero, which validates the pairing; the
probability-of-improvement family moves most.}
\label{tab:incumbentacq}
\begin{tabular}{lrrrr}
\toprule
acquisition & reference & treatment & $\Delta$ & share removed \\
\midrule
LogPI & 0.555 & 0.328 & $-$0.227 & +41\% \\
PI & 0.572 & 0.373 & $-$0.199 & +35\% \\
EI & 0.459 & 0.382 & $-$0.077 & +17\% \\
LogEI & 0.493 & 0.481 & $-$0.011 & +2\% \\
qNEI & 0.378 & 0.378 & +0.000 & +0\% \\
UCB & 0.346 & 0.346 & +0.000 & +0\% \\
qEI & 0.445 & 0.454 & +0.009 & $-$2\% \\
\bottomrule
\end{tabular}

\end{table}

\subsection{Breaking the descriptor confound}
\label{sec:manipulation}
Across any suite of real benchmarks $\optz$, sparsity and skew correlate at
about $0.9$ --- close to a structural identity for a bounded function with an
isolated optimum --- so Section~\ref{sec:mediation}'s descriptor terms cannot be
separated observationally, however many real benchmarks are added. Two
purpose-built families vary them by construction: a cos-field background plus one
spike, whose amplitude moves $\optz$ and whose width moves sparsity; and the Levy
function at $d \in \{4, 7, 11\}$, one landscape family across a dimension ladder
(Table~\ref{tab:manipulation}).

\textbf{$\optz$ on its own is not causal.} Holding the spike narrow and raising
its amplitude by $4\times$ moves $\optz$ from $5.0$ to $21.7$ and changes the
excess regret at $1\sigma$ from $0.246$ to $0.245$ --- no effect at all. Widen
the spike and the same amplitude change raises excess from $0.495$ to $1.888$.
What the observational $\beta_{\optz} = +0.42$ picks up is therefore an
interaction, not a main effect: a taller optimum costs more to miss only when it
is broad enough that the optimizer would otherwise have found it. On a true
needle, the optimizer was never going to find the optimum with or without noise,
so there is nothing for the noise to destroy. This also qualifies the reading of
the superlinear onset-0 behaviour offered in Section~\ref{sec:currency}: it is
not $\optz$ as such that makes early error expensive.

\textbf{Dimension, and a within-family ladder that corrects itself.} Within the
Levy family, excess regret at $1\sigma$ rises monotonically from $0.114$ at
$d = 4$ to $0.222$ at $d = 7$ and $0.264$ at $d = 11$, while the cross-suite
regression puts $\beta_d$ at $-0.03$ and not significant. Read on its own that
looks like a real effect the observational analysis missed. It is not, or not
mostly: $\optz$ also rises along the Levy ladder, $1.53$ to $2.07$ to $2.62$, so
that comparison varies dimension and signal strength together.

The \texttt{bump} ladder separates them. Amplitude and width are chosen at each
dimension so that $\optz$ is $9.0$ and the spike occupies a $10^{-6}$ fraction of
the box at all three, leaving dimension as the only difference. Excess regret is
then $0.031$ at $d = 4$, $0.028$ at $d = 7$ and $0.048$ at $d = 11$ --- not
monotone, and a factor of $1.5$ across the range where the confounded ladder gave
$2.3$. \textbf{Most of what looked like a dimension effect was $\optz$.} What
does rise monotonically with dimension is the inference metric, the true value of
the design a practitioner would actually deploy: $0.52$, $0.70$, $0.99$. Higher
dimension does not destroy more of the optimizer's progress; it makes the
progress harder to identify from noisy observations at the end.

The methodological point survives in a sharper form than we first wrote it. A
landscape-level regression over a hand-assembled suite can miss an effect, and a
within-family ladder is the cheap fix --- but a family indexed by dimension is not
automatically a manipulation of dimension, and ours was not until it was built to
be.

\begin{table}[t]
\centering
\caption{Purpose-built landscapes. Raw post-onset excess regret in landscape
standard deviations, gaussian error from iteration~1. Amplitude and width are
varied independently in the top block; the bottom block is one function family
at three dimensions.}
\label{tab:manipulation}
\begin{tabular}{lrrrrrrr}
\toprule
landscape & $d$ & $\mathrm{opt}_z$ & sparsity & $0.05\sigma$ & $0.25\sigma$ & $1\sigma$ & $5\sigma$ \\
\midrule
\texttt{bump\_a4\_w0.05} & 4 & 5.0 & $<$1.5e-05 & 0.007 & 0.152 & 0.246 & 0.324 \\
\texttt{bump\_a16\_w0.05} & 4 & 21.7 & $<$1.5e-05 & 0.014 & 0.026 & 0.245 & 0.262 \\
\addlinespace
\texttt{bump\_a4\_w0.15} & 4 & 5.3 & 3.8e-04 & 0.133 & 0.289 & 0.495 & 0.704 \\
\texttt{bump\_a16\_w0.15} & 4 & 12.3 & 1.8e-04 & 0.298 & 0.730 & 1.888 & 3.173 \\
\midrule
\texttt{levy\_4d} & 4 & 1.5 & 2.6e-01 & 0.006 & 0.048 & 0.114 & 0.200 \\
\texttt{levy\_7d} & 7 & 2.1 & 7.3e-02 & 0.051 & 0.115 & 0.222 & 0.305 \\
\texttt{levy\_10} & 11 & 2.6 & 1.3e-02 & 0.001 & 0.147 & 0.264 & 0.481 \\
\midrule
\texttt{bump\_d4} & 4 & 9.0 & $<$1.5e-05 & 0.001 & 0.121 & 0.279 & 0.327 \\
\texttt{bump\_d7} & 7 & 9.0 & $<$1.5e-05 & $-$0.001 & 0.103 & 0.248 & 0.428 \\
\texttt{bump\_d11} & 11 & 9.0 & $<$1.5e-05 & 0.142 & 0.215 & 0.427 & 0.734 \\
\bottomrule
\end{tabular}

\end{table}

\section{Does any of this hold for multiple objectives?}
\label{sec:multiobjective}

The applied setting is frequently multi-objective \citep{daulton2021parallel}
and everything above is scalar, so the generalisation cannot be assumed. We
repeat the design on seven \textsc{BoTorch} problems with a published maximum
hypervolume --- BraninCurrin, ZDT1, DTLZ2, DH2, VehicleSafety, Penicillin and
CarSideImpact --- spanning two to four objectives and two to seven dimensions,
with the four hypervolume acquisitions (qEHVI, qNEHVI and their log variants)
and the same two model-free floors, at five seeds: $1{,}680$ paired runs, of
which $960$ are model-based noisy-versus-clean comparisons on the problems that
survive the headroom screen below.

Standardisation transfers exactly rather than approximately. Each objective is
mapped to zero mean and unit variance over the same fixed Sobol sample used for
the scalar suite, and under a per-objective affine map the hypervolume scales by
$\prod_m \sigma_m$, so the standardised maximum follows in closed form from the
published one. We verify this rather than assert it: across the seven problems
the largest relative error between the closed-form standardised maximum and a
direct recomputation is $4.7 \times 10^{-16}$.

\paragraph{Controls.} The model-free floor is again identically zero across all
$560$ such comparisons, so no observation noise leaks into the candidate stream
in the hypervolume code path either. No run beats its published maximum
hypervolume. Acquisition-optimisation fallbacks total zero, which we report
rather than assume: a hypervolume acquisition that fails to construct falls back
to random sampling, and because that produces a plausible-looking curve rather
than an error, the arm has to be read off the fallback counter and not off the
plots.

\paragraph{One problem is dropped, by the same screen as before.} DH2's
reference point leaves a $32{,}768$-point Sobol sample with a hypervolume of
exactly zero, and its best learner does not improve on the model-free floor at
all. There is nothing for error to destroy, so its null is guaranteed by
construction; it is excluded. The remaining six pass, with Penicillin weakest at
$0.197$ --- comparable to Rosenbrock at $0.176$, the weakest admissible scalar
landscape.

\begin{table}[t]
\centering
\caption{The two arms on a common footing. Both columns are post-onset excess
divided by the gap between the optimum and a same-budget model-free floor, as a
percentage; brackets are 95\% cluster-bootstrap intervals over
landscapes/problems. This is \emph{not} the normalisation used in
Table~\ref{tab:dose}, and the two must not be read against each other.}
\label{tab:mo}
\begin{tabular}{lrrrr}
\toprule
& \multicolumn{2}{c}{scalar (20 landscapes)} & \multicolumn{2}{c}{multi-objective (6 problems)} \\
\cmidrule(lr){2-3} \cmidrule(lr){4-5}
error & \multicolumn{1}{c}{from trial 1} & \multicolumn{1}{c}{from trial 21} & \multicolumn{1}{c}{from trial 1} & \multicolumn{1}{c}{from trial 21} \\
\midrule
$0.05\sigma$ & 10.2 \makebox[3.4em][l]{\scriptsize [4, 19]} & 2.8 \makebox[3.4em][l]{\scriptsize [1, 5]} & 1.4 \makebox[3.4em][l]{\scriptsize [0, 3]} & 0.5 \makebox[3.4em][l]{\scriptsize [$-$1, 2]} \\
$0.25\sigma$ & 44.3 \makebox[3.4em][l]{\scriptsize [22, 71]} & 6.5 \makebox[3.4em][l]{\scriptsize [4, 10]} & 13.6 \makebox[3.4em][l]{\scriptsize [8, 19]} & 3.7 \makebox[3.4em][l]{\scriptsize [0, 10]} \\
$1\sigma$ & 99.0 \makebox[3.4em][l]{\scriptsize [54, 154]} & 10.3 \makebox[3.4em][l]{\scriptsize [7, 13]} & 49.1 \makebox[3.4em][l]{\scriptsize [42, 56]} & 10.7 \makebox[3.4em][l]{\scriptsize [2, 24]} \\
$5\sigma$ & 148.4 \makebox[3.4em][l]{\scriptsize [76, 239]} & 16.5 \makebox[3.4em][l]{\scriptsize [11, 22]} & 72.3 \makebox[3.4em][l]{\scriptsize [55, 91]} & 13.0 \makebox[3.4em][l]{\scriptsize [3, 27]} \\
\bottomrule
\end{tabular}

\end{table}

\paragraph{The comparison needs one piece of care.} Elsewhere the paper divides
excess regret by $\optz$, a z-score of the optimum against a random design's
mean. A Pareto problem has no scalar optimum and therefore no $\optz$, and the
natural hypervolume analogue --- the gap between the published maximum and what a
model-free design of the same budget reaches --- is a numerically much smaller
denominator. Quoting one arm's published fraction against the other's would
manufacture a threefold difference out of the choice of divisor alone. Both
columns of Table~\ref{tab:mo} are therefore recomputed on the floor denominator,
identically; and we also report the onset ratio, which is a ratio of two costs on
the same problem and so does not depend on the denominator at all.

\textbf{The two headline structures replicate.} Cost rises monotonically with
error magnitude and falls sharply when the error arrives late, on both arms. At
$1\sigma$ from the first rating the multi-objective arm loses $49.1\%$ of the
achievable hypervolume gain $[42, 56]$ against the scalar arm's $93.7\%$
$[52, 145]$; arriving at iteration~21 the two are indistinguishable, $10.7\%$
$[2, 24]$ against $10.7\%$ $[7, 14]$.

\textbf{There is no evidence multi-objective BO is the more fragile of the two,
and some that it is less.} The point estimates for early error are roughly half
the scalar arm's at every magnitude, and the intervals separate at $0.05\sigma$
and $0.25\sigma$; at $1\sigma$ they overlap slightly and at $5\sigma$
substantially, so the honest reading is a consistent direction rather than a
clean separation at the top of the range. The onset ratio is correspondingly
flatter: $3.0$--$5.5\times$ across magnitudes against the scalar arm's
$5.0$--$8.9\times$. A mechanism consistent with both observations, and with
Appendix~\ref{sec:incumbent}, is that a hypervolume objective has no single
incumbent to corrupt --- the state a noisy observation can damage is a Pareto
set, and an error large enough to promote one dominated point does not discard
the rest of the front.

\begin{table}[t]
\centering
\caption{Hypervolume acquisitions, ranked over problems within each condition.
Absolute loss is the fraction of the achievable hypervolume gain lost under
noise. The four are not distinguishable.}
\label{tab:moacq}
\begin{tabular}{lrr}
\toprule
acquisition & mean rank & absolute loss \\
\midrule
qNEHVI & 2.40 & 20.2\% \\
qLogEHVI & 2.44 & 20.4\% \\
qEHVI & 2.56 & 20.7\% \\
qLogNEHVI & 2.60 & 20.8\% \\
\bottomrule
\end{tabular}

\end{table}

\textbf{The choice among hypervolume acquisitions does not appear to matter.}
Mean ranks span $2.40$ to $2.60$ on a scale where $2.5$ is a tie, absolute losses
span $20.2\%$ to $20.8\%$, and the Friedman test reaches $p < 0.05$ in $0$ of
$8$ conditions (Table~\ref{tab:moacq}). Kendall's $W$ is $0.056$ --- against
$0.247$ on the scalar suite --- and each of the four acquisitions is the most
robust on at least one of the six problems. Six problems at five seeds is far
less power than twenty at ten, so the honest reading is ``no detectable
difference'' rather than ``no difference''. It does, though, point the same way as
Section~\ref{sec:acq}: the disagreement between design spaces about which
acquisition is most robust is, if anything, sharper here.

\paragraph{What does not transfer.} Equation~\ref{eq:frag} is defined for a
scalar objective: it is the loss from one greedy pick among $m$ candidates, and
on a Pareto problem there is no single greedy pick whose loss it measures. The
mediation result is therefore untested in this arm rather than confirmed or
refuted by it. A hypervolume analogue would have to define the one-shot loss over
hypervolume contributions and take a stand on how per-objective error propagates
into them; that is a paper of its own, and we do not claim
Table~\ref{tab:mediation} extends until someone writes it.

\section{Three design choices that could have produced the result}
\label{sec:designchoices}

Every number above was measured at a fixed budget, against an unbounded
continuous response, using standard acquisition functions. Each of those is a
choice, and each could be doing the work.

\begin{table}[t]
\centering
\caption{Budget, with the late onset held at $40\%$ of the run in each case so
that budget is the only thing that changes. Percentages are the fraction of the
achievable improvement destroyed.}
\label{tab:budget}
\begin{tabular}{llrrrr}
\toprule
budget & error present & $0.05\sigma$ & $0.25\sigma$ & $1\sigma$ & $5\sigma$ \\
\midrule
$T=25$ & from it.\ 1 & 0.6 & 3.2 & 9.4 & 15.4 \\
 & from it.\ 11 & 0.3 & 1.3 & 4.1 & 8.0 \\
 & ratio & 2.2$\times$ & 2.4$\times$ & 2.3$\times$ & 1.9$\times$ \\
\midrule
$T=50$ & from it.\ 1 & 0.6 & 3.8 & 11.7 & 21.3 \\
 & from it.\ 21 & 0.1 & 1.0 & 2.3 & 4.7 \\
 & ratio & 4.4$\times$ & 3.9$\times$ & 5.1$\times$ & 4.6$\times$ \\
\midrule
$T=100$ & from it.\ 1 & 0.2 & 3.4 & 11.1 & 23.9 \\
 & from it.\ 41 & 0.1 & 0.4 & 0.9 & 2.2 \\
 & ratio & 1.4$\times$ & 9.2$\times$ & 11.7$\times$ & 10.8$\times$ \\
\bottomrule
\end{tabular}

\end{table}

\paragraph{Budget, and the onset effect is not budget-free.} Re-running at $T=25$
and $T=100$ with the onset held at the same fraction of the run gives
Table~\ref{tab:budget}. Early-error cost is roughly flat in budget --- $9.4\%$,
$14.1\%$, $11.1\%$ at $1\sigma$ --- but late-error cost falls steeply, from
$4.1\%$ at $T=25$ to $2.8\%$ at $T=50$ to $0.9\%$ at $T=100$. The onset ratio
therefore grows with the budget, from $2.3\times$ to $5.0\times$ to $11.7\times$.
This is the clearest limitation on the headline: \textbf{the $5\times$ onset
effect is a statement about a fifty-iteration study.} The direction is
unsurprising, since more remaining iterations means more chance to recover, but
the size is not: a fourfold change in budget moves the ratio by five. We looked
for a rescaling that would let one budget speak for another and did not find
one. Late-onset cost per remaining iteration falls by $18\times$ across the three
budgets, and late-onset cost times remaining iterations is non-monotone
($0.62$, $0.84$, $0.54$), so neither is the invariant. The early-error number is
the transferable one; a study at a different budget has to measure its own onset
effect.

\begin{table}[t]
\centering
\caption{A bounded, discrete response. The instrument arm clips to $\pm 8\sigma$
and rounds to $0.55\sigma$, the range and step implied by the three archival
studies.}
\label{tab:instrument}
\begin{tabular}{llrrrr}
\toprule
response scale & error present & $0.05\sigma$ & $0.25\sigma$ & $1\sigma$ & $5\sigma$ \\
\midrule
unbounded, continuous & from it.\ 1 & 0.6 & 3.8 & 11.7 & 21.3 \\
 & from it.\ 21 & 0.1 & 1.0 & 2.3 & 4.7 \\
bounded, discrete & from it.\ 1 & 3.8 & 5.0 & 11.7 & 20.8 \\
 & from it.\ 21 & 1.3 & 1.4 & 2.5 & 4.6 \\
\bottomrule
\end{tabular}

\end{table}

\paragraph{A real rating scale, which cuts both ways.} The synthetic objective is
unbounded and continuous; a person returns a bounded, discrete rating. The three
archival instruments span $13$, $17$ and $22\sigma$ with steps of $0.55$, $0.54$
and $1.12\sigma$, so the arm clips to $\pm 8\sigma$ and rounds to $0.55\sigma$ ---
the real instrument, in the study's own currency. Table~\ref{tab:instrument}
shows the two effects separately. At $0.05\sigma$ the discretised arm is four
times worse ($3.8\%$ against $0.9\%$), because rounding to a $0.55\sigma$ grid is
itself an error of about $0.16\sigma$ and swamps an injected error far smaller
than the step. At $5\sigma$ it is slightly better ($20.8\%$ against $22.0\%$),
because clipping bounds how far one bad rating can move the incumbent. A rating
scale is a noise floor at the bottom of the range and a safety rail at the top.

\begin{table}[t]
\centering
\caption{Robust baselines. qKG is the knowledge gradient; replication spends
every second evaluation re-asking about the incumbent. Compared at a fixed number
of evaluations, so replication buys averaging with half the designs.}
\label{tab:robust}
\begin{tabular}{llrrrr}
\toprule
acquisitions & error present & $0.05\sigma$ & $0.25\sigma$ & $1\sigma$ & $5\sigma$ \\
\midrule
ten standard & from it.\ 1 & 1.1 & 5.6 & 13.9 & 21.5 \\
 & from it.\ 21 & 0.3 & 1.3 & 2.7 & 4.5 \\
qKG and replication & from it.\ 1 & 0.3 & 2.2 & 9.2 & 16.2 \\
 & from it.\ 21 & 0.2 & 0.1 & 2.2 & 4.8 \\
\bottomrule
\end{tabular}

\end{table}

\paragraph{Methods built for the problem recover about a third.} All twelve
acquisitions above are ordinary ones run under noise, which says which standard
choice survives best but not how much a purpose-built method would recover. Two
baselines: the knowledge gradient, which values the improvement in the posterior
\emph{maximum} rather than in any observed value, and the practical answer of
spending half the budget re-asking about the incumbent. At $1\sigma$ from the
first rating they lose $9.2\%$ of the achievable gain against the standard
acquisitions' $14.1\%$, and at $5\sigma$, $16.2\%$ against $22.0\%$
(Table~\ref{tab:robust}). At the late onset the two are close to
indistinguishable. So a third of the early-error cost is recoverable by choosing
a different method, at a fixed evaluation budget --- which for replication means
buying that recovery with half as many distinct designs. Two thirds is not
recoverable this way. Equation~\ref{eq:frag} measures that residue: the
information the error destroys, rather than any optimizer's failure to cope with
it.

\section{The claims replicate on seeds that had no part in making them}
\label{sec:confirmatory}

Everything above was selected and estimated on the same ten seeds. Before
running any fresh ones we wrote down four claims, the statistics that would test
them, the thresholds each had to clear, and four control conditions that would
void the run --- and then ran ten seeds that had taken no part in generating any
hypothesis.

\begin{table}[t]
\centering
\caption{Preregistered quantities on the screening seeds and on ten fresh ones.
The replication line reports the \emph{sensitivity} analysis described in the
text, not the preregistered run, which is void.}
\label{tab:confirmatory}
\begin{tabular}{lrr}
\toprule
preregistered quantity & screening & fresh seeds \\
\midrule
H1 $\beta(\frag)$ & +0.65 & +0.69 \\
H1 $\beta(\log_{10}\sigma_e)$ & $-$0.06 & $-$0.07 \\
H2 onset ratio & 5.44 & 5.01 \\
H3 robust mean rank & 4.24 & 4.11 \\
H3 fragile mean rank & 7.26 & 6.98 \\
H4 dose-response monotone & yes & yes \\
\midrule
\multicolumn{3}{l}{replicates: H1: yes; H2: yes; H3: yes; H4: yes} \\
\bottomrule
\end{tabular}

\end{table}

\textbf{The preregistered run is void, by its own rule.} One control failed:
Rosenbrock's headroom on the fresh seeds is $0.083$, below the $0.10$ the
preregistration fixed in advance. The document says that a control failure makes
the run void rather than negative, and an exclusion applied after seeing the data
cannot undo that. We report it as void.

We also report what it showed, labelled as the post-hoc sensitivity analysis it
is. Dropping Rosenbrock from both arms, all four claims clear their thresholds:
$\beta(\frag)$ moves from $+0.65$ to $+0.69$ while $\beta(\log\sigma_e)$ stays at
$-0.07$ with an interval containing zero; the onset ratio is $5.01$ against
$5.44$; the robust and fragile acquisition families sit at mean ranks $4.11$ and
$6.98$ against $4.24$ and $7.26$; and the dose-response is monotone in both
(Table~\ref{tab:confirmatory}).

One argument is worth making explicitly, because it decides how much the void
should worry a reader. The headroom control exists to stop a benchmark with
nothing to lose from contributing a null that looks like evidence of robustness.
It guards against false \emph{negatives}. All four claims came out positive, and
including a low-headroom benchmark makes a positive result harder to obtain, not
easier. The failure mode the control protects against is therefore not one that
can explain what we observed. That is a reason to read the sensitivity analysis
as informative; it is not a reason to call the run confirmed, and we do not.

\section{The landscape-descriptor fit in full}
\label{app:descriptors}

Section~\ref{sec:mediation} reports the descriptor model only as a comparison
for $\frag(\sigma_e)$. Table~\ref{tab:descriptors} gives it in full, with the
variance inflation factors that limit what can be read off it.

\begin{table}[h]
\centering
\caption{Post-onset excess regret on landscape descriptors and the error
magnitude, predictors $z$-scored, cluster-bootstrapped over landscapes with
Benjamini--Hochberg correction. VIF is computed on the design matrix.}
\label{tab:descriptors}
\begin{tabular}{lrcrr}
\toprule
term & $\beta$ & 95\% CI & $p_{\mathrm{FDR}}$ & VIF \\
\midrule
$\log\sigma_e$ & +0.37 & [+0.15, +0.66] & 0.001$^{**}$ & 1.0 \\
$\log\mathrm{opt}_z$ & +0.64 & [+0.24, +0.74] & 0.001$^{**}$ & 9.2 \\
$\log$ tail & +0.20 & [+0.03, +0.33] & 0.008$^{**}$ & 1.2 \\
rugged. & $-$0.09 & [$-$0.38, +0.12] & 0.519 & 1.3 \\
$d$ & $-$0.03 & [$-$0.26, +0.12] & 0.768 & 1.0 \\
\midrule
\multicolumn{5}{l}{$R^2 = 0.526$, 640 cells, 20 landscape clusters} \\
\bottomrule
\end{tabular}

\end{table}

$\optz$ carries a VIF of $9.2$ and regresses on the other descriptors with
$R^2 = 0.89$, so its coefficient and those of sparsity and skew (excluded here,
VIF $5.8$ and $6.3$) are not separately identified. This is the reason
Appendix~\ref{sec:manipulation} builds landscapes that vary these properties
independently rather than trying to disentangle them by regression. Dimension
and ruggedness are null, and Appendix~\ref{sec:manipulation} shows the null on
dimension is a property of the observational design rather than of dimension.

\section{Does noise ever help?}
\label{app:beneficial}

A noisy observation can in principle help, by pushing an over-exploiting
optimizer out of a local basin. At one seed per cell, $25$ of the $640$
landscape-by-condition cells showed reliably negative excess regret --- noise
apparently improving the outcome --- which is the kind of finding that invites a
mechanism.

At ten seeds it is one cell: \texttt{michalewicz} under \textsc{ar}(1) error at
$0.05\sigma$ from iteration $21$, mean excess $-0.136$, $d_z = -0.21$ over $100$
paired runs. One cell in $640$ at $p < 0.05$ is what a family of $640$ tests
produces when nothing is happening. We report it as a null and note the
sampling-variability lesson: the twenty-four cells that disappeared between one
seed and ten were all in the direction that would have made the better story.

\end{document}
